\input{uwhealthcv-pdfa}% PDF/A output (template/); must come before \documentclass
\documentclass[11pt]{article}

% The Faculty of Health format, APR highlighting and page-break rules all
% live in template/ (see template/README.md). This file holds content only.
% Compile main.tex for the plain Faculty of Health CV, or apr.tex for the
% APR version (evaluation notes + highlighting).
\usepackage{uwhealthcv}
\usepackage{pgfplots}% the HLTH 230 grades figure (APR version)
\pgfplotsset{compat=1.18}
% The aprletter environment, which the letters in letters/ are written in:
%
%   \begin{aprletter}{<recipient, lines separated by \\>}{<salutation>}
%     <text>
%   \aprclosing{Sincerely,}
%   \aprenclosure{<what is enclosed>}   (optional)
%   \end{aprletter}
%
% Read in the preamble of two kinds of document:
%  - a standalone letter (cover-letter.tex and the others: letter class, with
%    letterhead/format.tex), where it is the letter class's letter, \opening and \closing,
%    and the letterhead PDFs are laid over the first and last pages;
%  - main.tex, so that the APR version can include letters in the CV itself
%    (\ifAPR ... % Cover letter for the APR: the executive summary on SPHS letterhead, signed by Jim and
% placed on top of the APR packet (this letter, then apr.pdf). A separate document from
% the CV: compile it on its own (Overleaf: Menu > Main document > cover-letter.tex).
% Update the text each review cycle; counts that grow with the record are printed with
% \cvnumber (data/cv-numbers.tex, from scripts/generate.py).
%
% Sources for 2025--2026: Jim's prose draft (chat, 2026-11-24), revised with him so each
% section ends on recognition, and his correction that teaching includes supervision
% (so the course paragraph says "course teaching"). Figures from main.tex ("Significant
% Research Contributions", the Chickadee, HLTH 230, funding and supervision notes),
% data/venue-profile.tex (3 of the period's journal articles in health journals),
% courses.bib (all 2025--2026 offerings in 2026), grants.bib, service.bib (ih2026convener,
% ih2027tpc, cui2025gc, health2025council, health2025fusc, the appeals and SPHS
% committees), activities.bib (ih2026ecw, assessdht2025panel, nfrf2026exploration),
% scholar.bib (siek2027grand), awards.bib (gradsupervision2026nominee, cui2025service,
% chi2025reviews, dis2026reviews, rphs2025edir) and supervision.bib. Veen Wong defends on
% 30 November 2026 (Jim): once she has, say she completed her PhD.
%
% The text of the letter. The standalone document (cover-letter.tex) and the APR version
% of the CV (main.tex, \ifAPR) both read this file; letters/aprletter.tex defines the
% aprletter environment for each.
\begin{aprletter}{%
The Interim Director and the Annual Performance Review Committee\\
School of Public Health Sciences\\
University of Waterloo}{Dear Interim Director and Members of the SPHS APR Committee,}

Please find enclosed my annotated Faculty of Health CV for the 2025--2026 review period, which included a sabbatical from December 2024 to August 2025. I used the sabbatical to lay the groundwork for my lab's machine-learning research on qualitative analysis, which led directly to my students' recent work in this area (ICHEC 2025, CUI 2025, GI 2026). My research advances public health through human-centred computing in three programs: scalable tools for qualitative analysis of online health communities; co-designed games for health behaviour change and gender-norm shift; and technologies that help families care for aging parents at home. Across teaching, research and service, the period produced work that others now use and build on, and that peers have recognized.

\aprheading{research}{Research and Funding}

The work reached both health and computing audiences. Of my \cvnumber{review-journal} journal articles in the period, three appeared in health journals (two in \emph{ACM Transactions on Computing for Healthcare}, one in \emph{Social Science \& Medicine}), alongside peer-reviewed papers at ICHEC, CUI and Graphics Interface. Five of the six papers from my qualitative-analysis program (2022--25) are among the world's 10\% most cited for their field and year, led by Robert Gauthier's CHI 2022 paper in the top 2\% (SciVal). Bold Sky, a game co-designed with social work students to prevent intimate partner violence, was presented to practitioners at the International Summit on Ending Gender-Based Violence, and Veen Wong's presentation of it won the EDI-R Presentation Award at the 2025 Research and Practice in Health Sciences Conference. The National Research Council funded my aging-in-place program a second time (\$321K as Co-PI, 2026--28, after \$201K as PI). New awards also came from CIHR (\$100K, Co-applicant), the Graham Seed Fund (\$35K, PI) and the University, alongside my NSERC Discovery Grant, held continuously since 2015 and more than doubled at renewal. A \$1.1M CIHR Project Grant that I lead as Nominated Principal Applicant is under review.

\aprheading{teaching}{Course Teaching}

My course teaching in the period came in 2026, after my return from sabbatical, and centred on renewal; graduate supervision continued throughout. I redeveloped HLTH 230, and all six of its Student Course Perceptions teaching items rose across its two 2026 offerings. I also rebuilt HLTH 615 around human-centred design with the Centre for Extended Learning. For HLTH 230 I designed and built Chickadee, a programming and assessment platform that now runs as a production campus service. Hosted by IST, it uses single sign-on, sends grades to LEARN (an integration the Dean approved) and passed the University's Policy 46 risk and privacy assessments, with independent security testing. In its Spring 2026 pilot, 143 students made 8,904 submissions and got results in a median of 7.5 seconds. Chickadee is open source, SPHS colleagues plan to adopt it, and a Teaching Innovation Incubator proposal in preparation would bring it to more than 1,300 students in seven Health and Engineering courses in 2027--28. I present it at WatITis, the campus IT conference, on December 9.

\aprheading{teaching}{Graduate Supervision}

I funded nine of my ten graduate students from my grants, and trainees first-authored \cvnumber{review-works-trainee-first} of my \cvnumber{review-works} papers and presentations in the period. Three MMath students completed within program limits and moved into roles including Layer 6 (TD), and Veen Wong defends her PhD on November 30. I sat on 11 graduate committees in 2025 and 12 in 2026. My students nominated me for the University's Award of Excellence in Graduate Supervision.

\aprheading{service}{Service and Leadership}

I led two major international initiatives. As General Chair of ACM CUI 2025, I brought an international conference of about 130 attendees to Waterloo and received its Recognition of Service. As a founding convener of the ACM Conference on Interactive Health, I helped establish a new ACM venue for HCI and health research. Its first conference (Porto, July 2026) drew about 230 attendees in person; I co-chaired its Early Career Workshop for 15 early-career researchers and am its 2027 Technical Program Chair, overseeing peer review of its archival papers. I also co-organized ``Grand Challenges in HCI~\texttimes{} Health,'' a workshop accepted for CHI 2027. In governance, I served as Chair Elect (2025--26) and now Chair (2026--27) of the Health Faculty Council, and on the Faculty Undergraduate Studies Committee, two Faculty appeals committees and three SPHS program committees. I also served on the EU ASSESS-DHT expert panel on digital health technology assessment and an NFRF review panel, and my peer reviews received special recognition at CHI 2025 and DIS 2026.

The CV's notes place this record in the context of HCI's norms, with field-weighted citation evidence. Because the SPHS guidelines provide for expert opinion where assessing quality and impact calls for domain expertise, I ask that the Director obtain it from a University colleague with a regular appointment in HCI.

\aprclosing{Sincerely,}
\end{aprletter}
 ... \fi). There each letter starts a new page with
%    1-inch margins and no running header, the letterhead (300 dpi images of the same
%    PDFs, letterhead/*.png, so the CV stays PDF/A) behind its first and last pages, and
%    the CV's own layout and header resume after it.
\newcommand\aprsignature{Jim Wallace, PhD\\Professor\\School of Public Health Sciences\\University of Waterloo}
\makeatletter
\@ifclassloaded{letter}{%
  \signature{\aprsignature}%
  \date{}%
  \newenvironment{aprletter}[2]{\begin{letter}{#1}\opening{#2}}{\end{letter}}%
  \newcommand\aprclosing[1]{\closing{#1}}%
  \newcommand\aprenclosure[1]{\encl{#1}}%
}{%
  % #1: an image the size of the page, drawn behind the page being built
  \newcommand\apr@letterhead[1]{%
    \AddToHookNext{shipout/background}{%
      \put(0,0){\raisebox{-\paperheight}{%
        \ifcv@tagged\includegraphics[width=\paperwidth,height=\paperheight,artifact]{#1}%
        \else\includegraphics[width=\paperwidth,height=\paperheight]{#1}\fi}}}}%
  \newenvironment{aprletter}[2]{%
    \newgeometry{margin=1in}% starts a new page
    \pagestyle{empty}%
    \apr@letterhead{letterhead/SPHS_letterhead_page_1.png}%
    \setlength\parskip{0.6\baselineskip}%
    \par\noindent\hfill\usvardate\today\par\medskip
    \noindent#1\par\medskip
    \noindent#2\par}%
   {\par\restoregeometry}% ends the letter's last page; the group restores the page style
  \newcommand\aprclosing[1]{%
    \apr@letterhead{letterhead/SPHS_letterhead_page_2.png}%
    \par\nobreak\noindent#1\par\nobreak\vspace{2\baselineskip}%
    {\setlength\parskip{0pt}\noindent\aprsignature\par}}%
  \newcommand\aprenclosure[1]{\par\noindent Encl: #1\par}%
}
\makeatother
% the letters the APR version includes (letters/)

\addbibresource{scholar.bib}
\addbibresource{grants.bib}
\addbibresource{supervision.bib}
\addbibresource{awards.bib}
\addbibresource{service.bib}
\addbibresource{activities.bib}
\addbibresource{courses.bib}
\addbibresource{committees.bib}
% MASc and MMath committees count and list with MSc
\cvcommitteegroup{MSc}{MSc, MASc, MMath}
\cvowner{Jim Wallace}
\setreviewperiod{2025}{2026}
% Publications and presentations from the review period are highlighted as
% research. To use another colour for one entry, add its key to aprteaching,
% aprservice or apraward; to stop one being highlighted, add it to aprnone.




\DeclareBibliographyCategory{conference}
\addtocategory{conference}{haller2010nice,wallace2008contextual,wallace2008comparing,inkpen2008groupbanter,martinez2011modelling,martinez2011modellingb,kelly2010solar,wallace2013collaborative,wallace2014effect,wallace2014bezelsearch,pietroszek2014smartcasting,pietroszek2015tiltcasting,huijbregts2015talkingtiles,wallace2016creating,katsuragawa2016gestural,katsuragawa2016watchpoint,cheung2016felines,pietroszek2017watchcasting,wallace2017subtle,siddhpuria2017exploring,homaeian2018group,liu2018applying,ghare2018increasing,selwyn2019co,tahai2019scalebridge,bomfim2020food,gauthier2022will,Woo2024,kaur2025transformers,prasad2026pragmatica,wong2026ipvprevention,ugayamazza2026agency,sharma2026details}

% Extended abstracts, posters, demos, works in progress, interactivity, and the
% workshop and meetup proposals I co-organize (CHI EA, companion and adjunct
% proceedings): refereed, but not full or short papers, so they print apart.
\DeclareBibliographyCategory{abstracts}
\addtocategory{abstracts}{homaeian2017investigating, chan2018smartsurveys, chan2020post, chan2018changing, bomfim2018pirate, pietroszek20153d, ha2006my, wallacelightweight, CHI2025Workshop, scott2015local, wallaceexploring, wallace2006edit, wallace2006swordfishchiea, Wong2025DIS, AttentionGo, wong2025praise, avellino2026meetup, siek2027grand, avellino2027meetup, wallace2012exploring, wallace2012using, Sharma2025, cheung2016quantum, wallace2017disappearing, anslow2016collaboration, chen2019design}

\DeclareBibliographyCategory{journal}
\addtocategory{journal}{wallace2009investigating,inkpen2006left,wallace2011investigating,lee2016web,wallace2017technologies,aboueid2021young,aboueid2021use,homaeian2021joint,chan2021merlynne,reetz2021nature,aboueid2021latent,gauthier2022computational,mcmurray2018s,homaeian2022handoff,bomfim2022design,wada2022designing,gauthier2023agency,wei2023protocol,BARDWELL2024116683,homaeian2025,Wong2026IPV, Homaeian2026, homaeian2027participatory, rotolo2026investigating}

\DeclareBibliographyCategory{technicalreports}
\addtocategory{technicalreports}{wallace2026interdisciplinary,wallace2021roadmap,wallace2008interaction,kuzminykhchanging,amr}

\DeclareBibliographyCategory{workshops}
\addtocategory{workshops}{CHI2025Workshop,siek2027grand,wallace2017disappearing,scott2015local,anslow2016collaboration}



\DeclareBibliographyCategory{other}
\addtocategory{other}{wallace2006swordfish,wallace2012impact}

% \DeclareBibliographyCategory{thesis}
% \addtocategory{thesis}{wallace2006swordfish,wallace2012impact}

\DeclareBibliographyCategory{patent}
\addtocategory{patent}{quinn2012supporting}

\DeclareBibliographyCategory{software}
\addtocategory{software}{cta2022, curio, pass,wallace_chickadee,wallace2026waterloo,wallace2015quantum,wallace2026quantum}
% Releases, commits and lines of code under each software entry (data/software.yaml;
% regenerated by scripts/software_stats.py). \setentrydetail comes from the template
% (from the release after v0.2.0); until template/ has it, the lines don't print.
% Not printed from 24 November 2026: development volume is effort, not impact, and the
% Chickadee note gives its use instead (revision list Jim forwarded, section 4). The
% stats are still generated; uncomment the line to print them again.
% \ifdefined\setentrydetail% Generated by scripts/software_stats.py from data/software.yaml; do not edit.
% Printed under each entry in PUBLICATIONS > Software (\setentrydetail, template).
\setentrydetail{wallace_chickadee}{Developed February -- November 2026: 1,565 releases, 3,427 commits, 238,000 lines of code and 222,000 lines of automated tests.}
\setentrydetail{cta2022}{Developed April 2021 -- July 2023: 14 releases, 289 commits and 26,000 lines of code.}
\fi
% Citation impact under highly cited papers (SciVal export; regenerated by
% scripts/citation_impact.py).
\ifdefined\setentrydetail% Generated by scripts/citation_impact.py from data/citation-impact.yaml; do not edit.
% Entries in the top 10% of field-weighted citations (SciVal), published
% before 2026; printed under the entry in both PDFs (\setentrydetail, template).
\setentrydetail{kaur2025transformers}{\textbf{Top 1\% most cited in its field} (field-weighted citation impact 20.99; SciVal, November 2026, indexed as 2026)}
\setentrytag{kaur2025transformers}{Top 1\%}
\setentrydetail{gauthier2022will}{\textbf{Top 2\% most cited in its field} (field-weighted citation impact 5.84; SciVal, November 2026)}
\setentrytag{gauthier2022will}{Top 2\%}
\setentrydetail{haller2010nice}{\textbf{Top 3\% most cited in its field} (field-weighted citation impact 5.93; SciVal, November 2026)}
\setentrytag{haller2010nice}{Top 3\%}
\setentrydetail{wallace2013collaborative}{\textbf{Top 3\% most cited in its field} (field-weighted citation impact 5.08; SciVal, November 2026)}
\setentrytag{wallace2013collaborative}{Top 3\%}
\setentrydetail{pietroszek2015tiltcasting}{\textbf{Top 6\% most cited in its field} (field-weighted citation impact 3.34; SciVal, November 2026)}
\setentrytag{pietroszek2015tiltcasting}{Top 6\%}
\setentrydetail{Sharma2025}{\textbf{Top 7\% most cited in its field} (field-weighted citation impact 3.31; SciVal, November 2026)}
\setentrytag{Sharma2025}{Top 7\%}
\setentrydetail{wallace2009investigating}{\textbf{Top 7\% most cited in its field} (field-weighted citation impact 3.26; SciVal, November 2026)}
\setentrytag{wallace2009investigating}{Top 7\%}
\setentrydetail{BARDWELL2024116683}{\textbf{Top 7\% most cited in its field} (field-weighted citation impact 3.18; SciVal, November 2026)}
\setentrytag{BARDWELL2024116683}{Top 7\%}
\setentrydetail{gauthier2022computational}{\textbf{Top 7\% most cited in its field} (field-weighted citation impact 3.12; SciVal, November 2026)}
\setentrytag{gauthier2022computational}{Top 7\%}
\setentrydetail{homaeian2025}{\textbf{Top 9\% most cited in its field} (field-weighted citation impact 2.86; SciVal, November 2026)}
\setentrytag{homaeian2025}{Top 9\%}
\setentrydetail{martinez2011modelling}{\textbf{Top 9\% most cited in its field} (field-weighted citation impact 2.77; SciVal, November 2026)}
\setentrytag{martinez2011modelling}{Top 9\%}
\setentrydetail{katsuragawa2016watchpoint}{\textbf{Top 9\% most cited in its field} (field-weighted citation impact 2.62; SciVal, November 2026)}
\setentrytag{katsuragawa2016watchpoint}{Top 9\%}
\setentrydetail{wallace2014bezelsearch}{\textbf{Top 10\% most cited in its field} (field-weighted citation impact 2.46; SciVal, November 2026)}
\setentrytag{wallace2014bezelsearch}{Top 10\%}
\fi
% What each entry is missing, marked in orange only in the review copy (gaps.tex;
% regenerated by scripts/generate.py).
% Generated by scripts/gaps.py from the .bib files, main.tex, data/metrics.yaml and
% data/gaps.yaml; do not edit. Printed only in the review copy (gaps.tex). The
% definitions below stand in until the template provides \setentrygap and \cvgap.
\providecommand\setentrygap[2]{}
\providecommand\cvgap[1]{}
% 47 entries with gaps
\setentrygap{kelly2010solar}{acceptance rate or impact factor (unresolved, last looked for 2026-11-18)}
\setentrygap{wallace2008contextual}{acceptance rate or impact factor (unresolved, last looked for 2026-11-18)}
\setentrygap{wallaceexploring}{pages}
\setentrygap{wallacelightweight}{DOI, pages}
\setentrygap{hawant}{DOI, pages}
\setentrygap{pietroszek2017watchcasting}{acceptance rate or impact factor (unresolved, last looked for 2026-11-18)}
\setentrygap{tahai2019scalebridge}{acceptance rate or impact factor (unresolved, last looked for 2026-11-18)}
\setentrygap{homaeian2021joint}{acceptance rate or impact factor (unresolved, last looked for 2026-11-18)}
\setentrygap{chen2019design}{pages}
\setentrygap{wallace2006edit}{DOI}
\setentrygap{homaeian2022handoff}{acceptance rate or impact factor (unresolved, last looked for 2026-11-18)}
\setentrygap{bomfim2022design}{acceptance rate or impact factor (unresolved, last looked for 2026-11-18)}
\setentrygap{wada2022designing}{acceptance rate or impact factor (unresolved, last looked for 2026-11-18)}
\setentrygap{gauthier2023agency}{acceptance rate or impact factor (unresolved, last looked for 2026-11-18)}
\setentrygap{Wong2026IPV}{final publication details (in press), pages}
\setentrygap{Homaeian2026}{final publication details (in press), pages}
\setentrygap{homaeian2027participatory}{final publication details (in press), DOI, pages}
\setentrygap{cmentowski2024pdf}{start year}
\setentrygap{homaeianpdf}{start year}
\setentrygap{pietroszek2015phd}{start year}
\setentrygap{bomfim2022phd}{start year}
\setentrygap{homaeian2022phd}{start year}
\setentrygap{gauthier2022phd}{start year}
\setentrygap{wong2026phd}{start year, position}
\setentrygap{liphd}{start year}
\setentrygap{liaqatphd}{start year}
\setentrygap{ojiphd}{start year}
\setentrygap{chenphd}{start year}
\setentrygap{ng2017msc}{program, start year}
\setentrygap{tahai2018msc}{program, start year}
\setentrygap{chan2019msc}{program, start year}
\setentrygap{wada2019msc}{program, start year}
\setentrygap{woo2022msc}{program, start year}
\setentrygap{ugayamazza2023msc}{program, start year}
\setentrygap{kaur2024mmath}{program, start year}
\setentrygap{sharma2025mmath}{program, start year}
\setentrygap{prasad2026mmath}{program, start year}
\setentrygap{zhelizniak2026mmath}{program, start year}
\setentrygap{bradburymmath}{program, start year}
\setentrygap{leemmath}{program, start year}
\setentrygap{mainville2016mhi}{program, start year}
\setentrygap{mcnab2024msw}{program, start year, position}
\setentrygap{wallacejarvis2024msw}{program, start year, position}
\setentrygap{hlth454w2024}{SCP responses, SCP scores}
\setentrygap{hlth615f2020}{SCP responses, SCP scores}
\setentrygap{hlth454w2020}{SCP responses, SCP scores}
\setentrygap{hlth230w2020}{SCP responses, SCP scores}

% Counts the APR narratives quote, printed with \cvnumber (regenerated by
% scripts/generate.py; python3 scripts/cv_numbers.py --list names them).
% Generated by scripts/cv_numbers.py from scholar.bib, main.tex, grants.bib and data/*.yaml;
% do not edit. main.tex prints a number with \cvnumber{name}; the definitions below
% stand in until the template provides \setcvnumber and \cvnumber.
\providecommand\setcvnumber[2]{\expandafter\def\csname cvnumber/#1\endcsname{#2}}
\providecommand\cvnumber[1]{\ifcsname cvnumber/#1\endcsname\csname cvnumber/#1\endcsname\else\textbf{??}\fi}
\setcvnumber{review-journal}{5}
\setcvnumber{review-journal-trainee-first}{4}
\setcvnumber{review-conference}{13}
\setcvnumber{review-conference-trainee-first}{9}
\setcvnumber{review-presentations}{18}
\setcvnumber{review-presentations-trainee-first}{15}
\setcvnumber{review-works}{36}
\setcvnumber{review-works-trainee-first}{28}
\setcvnumber{review-fullpapers}{5}
\setcvnumber{review-fullpapers-trainee-first}{5}
\setcvnumber{review-abstracts}{8}
\setcvnumber{review-abstracts-trainee-first}{4}
\setcvnumber{scival-last-year}{2025}
\setcvnumber{scival-papers}{64}
\setcvnumber{scival-top}{13}
\setcvnumber{scival-top-share}{20}
\setcvnumber{scival-top-trainee-first}{6}
\setcvnumber{journals}{24}
\setcvnumber{journals-health}{9}
\setcvnumber{proceedings}{60}
\setcvnumber{proceedings-health}{2}
\setcvnumber{grants-funded}{34}
\setcvnumber{grants-funded-health}{6}

% Where work has been deployed or used beyond the paper (hand-maintained, with a
% source for each line; see the header of data/impact.tex).
\ifdefined\setentrydetail% Deployment and outreach impact: where work has been used beyond its paper,
% printed under the entry in both PDFs (\setentrydetail, from the template).
%
% Hand-maintained. Every fact needs a source you can point to: a public page from
% the organization that deployed it (IQC, a museum, a school board) or something
% Jim wrote in a cv-item issue or confirmed by email. Give the source in a
% "% Source:" comment above the line. No reach numbers (players, downloads,
% schools) without a source that states them.
%
% "% Search:" lists the names the weekly job searches for in Outlook and on the
% web (automation/weekly-update.md, step 2i).
%
% One line per key, and never for a key that data/software-stats.tex or
% data/citation-impact.tex already gives a line (\setentrydetail keeps only the
% last one); scripts/check_cv.py checks this.

% Search: "Quantum Cats", "Quantum Arcade"
% Source: Games Institute, "Games at the IQC Open House" (2015-10-02);
%   IQC, QUANTUM: The Exhibition page (THEMUSEUM programme, November 2016);
%   Games Institute, "Quantum Cats at the Ontario Science Centre" (2018-11-19);
%   IQC teacher resources page and Quantum Cats Educator's Guide (PDF, 2025).
%   At every stop of the exhibition: Jim (chat, 2026-11-22); venues and dates from
%   IQC's exhibition page.
\setentrydetail{wallace2015quantum}{Played at the IQC Open House (October 2015). Featured at all seven stops of IQC's touring \emph{QUANTUM: The Exhibition}: THEMUSEUM, Kitchener (2016); Science World, Vancouver; Western Development Museum, Saskatoon; TELUS Spark, Calgary; Discovery Centre, Halifax (2017); Canada Science and Technology Museum, Ottawa (2017--18); and Ontario Science Centre, Toronto (2018--19). An IQC classroom resource with an Educator's Guide (2025).}
\fi
% Publication metrics in the printed entries (Jim, 2026-11-25): Journal Impact
% Factors never print, and an acceptance rate prints only when it is under 30%.
% scholar.bib keeps every figure (the record, which scripts/gaps.py and the
% weekly run read); these maps only drop them from the printed note. To print
% one again, delete its step.
\DeclareSourcemap{
  \maps[datatype=bibtex]{
    \map[overwrite]{
      \step[fieldsource=note,
            match=\regexp{(\A|\.\s*)Journal\s+Impact\s+Factor\s*\(\d{4}\):\s*[\d.]+},
            replace={}]
      \step[fieldsource=note,
            match=\regexp{(\A|\.\s*)Acceptance\s+rate:\s*(\$\\sim\$\s*)?([3-9]\d|\d{3,})(\.\d+)?\\\%(\s*\([^)]*\))?},
            replace={}]
      \step[fieldsource=note, match=\regexp{\A\.\s*}, replace={}]
    }
    \map{
      \step[fieldsource=note, match=\regexp{\A\s*\Z}, final]
      \step[fieldset=note, null]
    }
  }
}


\DeclareBibliographyCategory{presentations}
\addtocategory{presentations}{Kaur2024cpi,Kaur2024,wallace2022activated,wallace2006interfacing,Homaeian2024,Wong2024,gauthier2022reflections,ng2017smartsurveys,agewell2023,gauthier2019understanding, wada2019rheumatoid, homaeian2019visualizing, chan2019leveraging, homaeian2019group,levy2019building,Chan2018,bomfim2019engaging, chan2018lifehack, wada2018evaluation, huijbregts2019effectiveness,chiplay2019foodworkshop,reconsidering2023,wish2023,homaeian2016designing,homaeianCHIWS2025,oji2025advancing,sharma2025thematic,wong2025SPHS,homaeian2025boundaries,homaeian2025community,ramirez2025prompt,wong2025boldsky,wong2025bold,joi2026privacyAAL,li2026glimpse,homaeian2026homelink,zhelizniak2026llm,wong2026IH,homaeian2026IH,wallace2026watitis, Kuzminykh2014, chan2018changingsymposium, kaur2024comparison, sharma2025details, prasad2025pragmatica}

% Media coverage and news releases about my work (published pieces only).
\DeclareBibliographyCategory{media}
\addtocategory{media}{uwnews2021symptomcheckers}

% Kept in scholar.bib but not printed (Jim's decisions). The SPHS salary analysis is
% listed once, as the technical report, which links its data and code (Jim, 2026-11-24).
\DeclareBibliographyCategory{unprinted}
\addtocategory{unprinted}{wallace2026sphssalary, hawant, termite_data_server, termite_ui, tomodapi}

\DeclareBibliographyCategory{invited}
\addtocategory{invited}
{hcischool2025,delpapa2025vibe}

% APR version: entries that also count as teaching and supervision, shown in
% both colours. Rule: a trainee I supervise is an author or presenter.
% Add new entries here as they appear.
\addtocategory{apralsoteaching}{Sharma2025, sharma2025details, prasad2025pragmatica, Wong2025DIS, hcischool2025, homaeian2025, homaeian2025boundaries, homaeian2025community, homaeianCHIWS2025, kaur2025transformers, oji2025advancing, ramirez2025prompt, sharma2025thematic, wong2025SPHS, wong2025bold, wong2025boldsky, AttentionGo, Homaeian2026, Wong2026IPV, prasad2026pragmatica, sharma2026details, ugayamazza2026agency, wong2026ipvprevention, li2026glimpse, homaeian2026homelink, zhelizniak2026llm, wong2026IH, homaeian2026IH, homaeian2027participatory, rphs2025edir, wong2025praise}
% Accepted during the review period but dated after it (published 2027).
\addtocategory{aprinclude}{homaeian2027participatory, siek2027grand, avellino2027meetup}
% Grant applications submitted in the review period and awaiting a decision (grants.bib).
\addtocategory{aprinclude}{cihr2026fallproject, sshrc2026insight, digitalscience2026catalyst, anthropic2026wellbeing, nserc2026create, mitacs2026adhd}
% Trainee work in the period that does not count as teaching (check_cv.py only; the
% template ignores it): the 2026 browser version of Quantum Cats credits its 2015 team,
% and Nast's co-op term was ten years ago (Jim, chat, 2026-11-22).
\DeclareBibliographyCategory{aprnoteaching}
\addtocategory{aprnoteaching}{wallace2026quantum}
% homaeian2025 (as published) and the CHI 2025 review recognitions were also in
% the 2023--24 APR, but Jim counts them in this review too (chat, 2026-11-20),
% so neither goes in aprnone.
% Chickadee is teaching software first, with research as its second colour.
\addtocategory{aprteaching}{wallace_chickadee, wallace2026watitis}
% The Faculty of Health CV template is service to colleagues first, with research second.
% So is the SPHS salary analysis report: service first, research second (#90).
\addtocategory{aprservice}{wallace2026waterloo, wallace2026interdisciplinary}
\addtocategory{apralsoresearch}{wallace2026waterloo, wallace2026interdisciplinary}
\addtocategory{apralsoresearch}{wallace_chickadee, wallace2026watitis, ih2027tpc, ih2026convener, ih2026ecw, cui2025gc, chiplay2023gc, iss2022pc}
% Reviewing, program committees, editorial roles and review or expert panels (activities.bib
% levels editorial, panel, grantreferee, journal, programcommittee, referee) are also evidence
% of scholarship (Faculty of Health guidelines). Add new entries of those levels here.
\addtocategory{apralsoresearch}{assessdht2025panel, nfrf2026exploration, nfrf2026international, nserc2015discoveryreview, nserc2025alliance, nserc2019polyani, nserc2019create, mitacs2013elevate, pacmgroup2024ae, tch2025, tweb2024, pacmiss2022editor, pacmchiplay2020ae, pacmcscw2019ae, pervasive2015, chb2015, cga2012, ijhcs2011, chi2018pc, cscw2019pc, gi2015pc, iss2024pc, perdis2015pc, chiplay2014pc, ppd2012pc, hypertext2010pc, chi2009referee, iss2009referee, cscw2008referee, uist2016referee, mobilehci2012referee, sui2013referee, mum2013referee, gi2016referee, iui2012referee, eics2012referee, dis2010referee, chiplay2014referee}
% Special recognitions for exceptional reviews, and the CHI PLAY 2023 and CUI 2025
% recognitions of service, are awards for service work.
% Chairing a student's examination: a committee role and service (committees.bib).
\addtocategory{apralsoservice}{hassan2026chair, videnov2023chair, lake2022chair, norman2019chair, mercer2020chair, sahu2018committee, payandeh2026examiner, xie2023examiner, yeh2023examiner, wehbe2021examiner, schaekermann2020examiner, surale2020examiner, tondello2019examiner}
\addtocategory{apralsoservice}{dis2026reviews, chi2025reviews, chi2023reviews, chi2022reviews, chi2020reviews, chi2017reviews, dis2016review, chi2015review, chiplay2014review, chiplay2023service, cui2025service}
% Workshops and meetups I co-organize at conferences print with the conference papers
% (refereed, with a DOI). A workshop proposal is scholarship first, with service second;
% a meetup is service to the community first, with teaching (mentorship) and research
% (two apralso... lists give a third colour, which the dots style shows).
\addtocategory{apralsoservice}{CHI2025Workshop, siek2027grand}
\addtocategory{aprservice}{avellino2026meetup, avellino2027meetup}
\addtocategory{apralsoteaching}{avellino2026meetup, avellino2027meetup}
\addtocategory{apralsoresearch}{avellino2026meetup, avellino2027meetup}


\begin{document}

\ifAPR % APR version only: the cover letter (the executive summary) on SPHS letterhead,
% before the CV; its text is letters/cover.tex (also cover-letter.tex on its own).
% Cover letter for the APR: the executive summary on SPHS letterhead, signed by Jim and
% placed on top of the APR packet (this letter, then apr.pdf). A separate document from
% the CV: compile it on its own (Overleaf: Menu > Main document > cover-letter.tex).
% Update the text each review cycle; counts that grow with the record are printed with
% \cvnumber (data/cv-numbers.tex, from scripts/generate.py).
%
% Sources for 2025--2026: Jim's prose draft (chat, 2026-11-24), revised with him so each
% section ends on recognition, and his correction that teaching includes supervision
% (so the course paragraph says "course teaching"). Figures from main.tex ("Significant
% Research Contributions", the Chickadee, HLTH 230, funding and supervision notes),
% data/venue-profile.tex (3 of the period's journal articles in health journals),
% courses.bib (all 2025--2026 offerings in 2026), grants.bib, service.bib (ih2026convener,
% ih2027tpc, cui2025gc, health2025council, health2025fusc, the appeals and SPHS
% committees), activities.bib (ih2026ecw, assessdht2025panel, nfrf2026exploration),
% scholar.bib (siek2027grand), awards.bib (gradsupervision2026nominee, cui2025service,
% chi2025reviews, dis2026reviews, rphs2025edir) and supervision.bib. Veen Wong defends on
% 30 November 2026 (Jim): once she has, say she completed her PhD.
%
% The text of the letter. The standalone document (cover-letter.tex) and the APR version
% of the CV (main.tex, \ifAPR) both read this file; letters/aprletter.tex defines the
% aprletter environment for each.
\begin{aprletter}{%
The Interim Director and the Annual Performance Review Committee\\
School of Public Health Sciences\\
University of Waterloo}{Dear Interim Director and Members of the SPHS APR Committee,}

Please find enclosed my annotated Faculty of Health CV for the 2025--2026 review period, which included a sabbatical from December 2024 to August 2025. I used the sabbatical to lay the groundwork for my lab's machine-learning research on qualitative analysis, which led directly to my students' recent work in this area (ICHEC 2025, CUI 2025, GI 2026). My research advances public health through human-centred computing in three programs: scalable tools for qualitative analysis of online health communities; co-designed games for health behaviour change and gender-norm shift; and technologies that help families care for aging parents at home. Across teaching, research and service, the period produced work that others now use and build on, and that peers have recognized.

\aprheading{research}{Research and Funding}

The work reached both health and computing audiences. Of my \cvnumber{review-journal} journal articles in the period, three appeared in health journals (two in \emph{ACM Transactions on Computing for Healthcare}, one in \emph{Social Science \& Medicine}), alongside peer-reviewed papers at ICHEC, CUI and Graphics Interface. Five of the six papers from my qualitative-analysis program (2022--25) are among the world's 10\% most cited for their field and year, led by Robert Gauthier's CHI 2022 paper in the top 2\% (SciVal). Bold Sky, a game co-designed with social work students to prevent intimate partner violence, was presented to practitioners at the International Summit on Ending Gender-Based Violence, and Veen Wong's presentation of it won the EDI-R Presentation Award at the 2025 Research and Practice in Health Sciences Conference. The National Research Council funded my aging-in-place program a second time (\$321K as Co-PI, 2026--28, after \$201K as PI). New awards also came from CIHR (\$100K, Co-applicant), the Graham Seed Fund (\$35K, PI) and the University, alongside my NSERC Discovery Grant, held continuously since 2015 and more than doubled at renewal. A \$1.1M CIHR Project Grant that I lead as Nominated Principal Applicant is under review.

\aprheading{teaching}{Course Teaching}

My course teaching in the period came in 2026, after my return from sabbatical, and centred on renewal; graduate supervision continued throughout. I redeveloped HLTH 230, and all six of its Student Course Perceptions teaching items rose across its two 2026 offerings. I also rebuilt HLTH 615 around human-centred design with the Centre for Extended Learning. For HLTH 230 I designed and built Chickadee, a programming and assessment platform that now runs as a production campus service. Hosted by IST, it uses single sign-on, sends grades to LEARN (an integration the Dean approved) and passed the University's Policy 46 risk and privacy assessments, with independent security testing. In its Spring 2026 pilot, 143 students made 8,904 submissions and got results in a median of 7.5 seconds. Chickadee is open source, SPHS colleagues plan to adopt it, and a Teaching Innovation Incubator proposal in preparation would bring it to more than 1,300 students in seven Health and Engineering courses in 2027--28. I present it at WatITis, the campus IT conference, on December 9.

\aprheading{teaching}{Graduate Supervision}

I funded nine of my ten graduate students from my grants, and trainees first-authored \cvnumber{review-works-trainee-first} of my \cvnumber{review-works} papers and presentations in the period. Three MMath students completed within program limits and moved into roles including Layer 6 (TD), and Veen Wong defends her PhD on November 30. I sat on 11 graduate committees in 2025 and 12 in 2026. My students nominated me for the University's Award of Excellence in Graduate Supervision.

\aprheading{service}{Service and Leadership}

I led two major international initiatives. As General Chair of ACM CUI 2025, I brought an international conference of about 130 attendees to Waterloo and received its Recognition of Service. As a founding convener of the ACM Conference on Interactive Health, I helped establish a new ACM venue for HCI and health research. Its first conference (Porto, July 2026) drew about 230 attendees in person; I co-chaired its Early Career Workshop for 15 early-career researchers and am its 2027 Technical Program Chair, overseeing peer review of its archival papers. I also co-organized ``Grand Challenges in HCI~\texttimes{} Health,'' a workshop accepted for CHI 2027. In governance, I served as Chair Elect (2025--26) and now Chair (2026--27) of the Health Faculty Council, and on the Faculty Undergraduate Studies Committee, two Faculty appeals committees and three SPHS program committees. I also served on the EU ASSESS-DHT expert panel on digital health technology assessment and an NFRF review panel, and my peer reviews received special recognition at CHI 2025 and DIS 2026.

The CV's notes place this record in the context of HCI's norms, with field-weighted citation evidence. Because the SPHS guidelines provide for expert opinion where assessing quality and impact calls for domain expertise, I ask that the Director obtain it from a University colleague with a regular appointment in HCI.

\aprclosing{Sincerely,}
\end{aprletter}

\setcounter{page}{1}% the CV's pages are numbered from its own first page
\fi

\noindent

\thispagestyle{empty}
\begin{center}
%\textbf{\textit{FACULTY OF HEALTH VITAE FORMAT}} \\
\cvsection{CURRICULUM VITA}
\end{center}

\cvsection{NAME/DATE:} \cvname, \usvardate\today
\\ \\
\cvsection{CONTACT INFORMATION:} \\
School of Public Health Sciences \\
University of Waterloo \\
james.wallace@uwaterloo.ca \\
ORCID: \extlink{https://orcid.org/0000-0002-5162-0256}{0000-0002-5162-0256} \\
Google Scholar: \extlink{https://scholar.google.ca/citations?user=uCxDo9sAAAAJ}{scholar.google.ca/citations?user=uCxDo9sAAAAJ} \\
Lab: \extlink{https://uwaterloo.ca/human-computer-interaction-health-lab/}{uwaterloo.ca/human-computer-interaction-health-lab}
\\ \\
\cvsection{DEGREES RECEIVED:} 

\begin{table}[h!]
\centering
\begin{tabular}{p{2.5cm} p{5cm} p{5cm}}
\underline{Degree}              & \underline{Institution}              & \underline{Years} \\
PhD                 & University of Waterloo   & 2007 – 2012  \\
MCSc                & Dalhousie University     & 2004 – 2006  \\
BMath               & University of Waterloo   & 1999 – 2004  \\
\end{tabular}
\end{table}

\cvsection{EMPLOYMENT HISTORY:}
\begin{table}[h!]
\centering
\begin{tabular}{p{3.5cm} l l}
\underline{Date}                & \underline{Position}                 & \underline{Institution} \\
2025 -- Present     & Professor                & University of Waterloo \\
2025 -- Present     & Professor (cross)        & Cheriton School of Computer Science \\
2020 -- 2025        & Associate Professor (cross) & Cheriton School of Computer Science \\
2019 – 2025         & Associate Professor      & University of Waterloo  \\
2014 – 2019         & Assistant Professor      & University of Waterloo   \\
2013 – 2014         & Assistant Professor      & Wilfrid Laurier University \\
2012 – 2013         & Post-Doctoral Researcher & University of Waterloo  \\
2010                & Intern                   & Google     \\
2006                & Intern                   & Microsoft Research 
\end{tabular}
\end{table}

\sectionbreak
\cvsection{ACADEMIC AWARDS AND DISTINCTIONS:}

% awards.bib: each @award has a level that picks its list; trainees' awards
% print under GRADUATE STUDENT SUPERVISION (\printawards[trainee]).
\cvsubheading{Research Awards:}
    \printawards[research]

\cvsubheading{Teaching and Supervision:}
    \printawards[teaching]

\cvsubheading{Recognition for Peer Review and Service:}
    \printawards[recognition]


\sectionbreak
\cvsection{SCHOLARLY AND PROFESSIONAL ACTIVITIES:}
\begin{cvlist}[0pt]
    \cvitem ACM Professional Member 2008 -- Present
\end{cvlist}


\cvsubheading{Editorial Responsibilities:}
    \printservice{editorial}

\cvsubheading{Review Panels and Expert Committees:}
    \printservice{panel}

\cvsubheading{External Tenure, Promotion and Program Reviews:}
    \printservice{externalreview}

\cvsubheading{Grant Referee:}
    \printservice{grantreferee}

\cvsubheading{Journal Referee:}
    \printservice{journal}

\cvsubheading{Conference Organizing Committee:}
    \printservice{organizing}

\cvsubheading{Conference Program Committee:}
    \printservice{programcommittee}

\cvsubheading{Conference Referee:}
    \printservice{referee}

\ifAPR % APR version only ------------------------------------------
%\newpage
\sectionbreak
\cvsection{A Note on Evaluating Human-Computer Interaction Research}
\\ \\
Human–Computer Interaction (HCI) research develops and evaluates interactive technologies, methods and conceptual frameworks that improve human outcomes, including health. It combines technological innovation with human-centred design and empirical evaluation, examining how people experience and use new systems in practice. In health research, HCI contributes the technologies, interaction methods and design knowledge through which interventions and health services can be developed and improved. It complements clinical and behavioural research by addressing how people engage with these interventions in practice (Blandford et al., 2018, \emph{Digital Health} 4). The field's leading archival publications appear in selective international conferences of the Association for Computing Machinery (ACM), which serve a role comparable to leading journals in other disciplines. Many now publish in the Proceedings of the ACM (PACM) series as journal articles. Full papers are reviewed anonymously by at least three experts, overseen by associate chairs and decided by an international program committee, and the Computing Research Association's 2025 guidance places such publication on par with or ahead of journals. 
\\ \\
\textbf{Paper awards.} Research excellence in HCI is also recognized through highly selective conference awards. At CHI (the field's largest venue), awards go to the top 5\% of submissions (Best Paper and Honourable Mention); the top 1\% receive Best Paper \extlink{https://chi2025.acm.org/best-paper-and-honorable-mention-awards-at-chi-2025/}{(CHI 2025 Awards)}. My CHI 2020 Honourable Mention was among the top 5\% of 3{,}126 submissions. The CSCW 2018 Honourable Mention was one of 19 that year \extlink{https://medium.com/acm-cscw/announcing-the-best-of-cscw-2018-b98cb91e0f61}{(Best of CSCW 2018)}.
\\ \\
\textbf{Authorship.} HCI research is collaborative, bringing together researchers with complementary technical, disciplinary and methodological expertise. \textbf{Graduate students frequently lead projects as first authors, while supervising researchers commonly provide intellectual and research leadership as last authors.} ACM's authorship policy requires substantial intellectual contributions and accountability for published work.
\\ \\
\textbf{Reading the highlighting.} Publications and presentations involving students I supervise (as authors, marked *, or presenters, marked \dag) represent both scholarly contributions and outcomes of graduate training. They are therefore highlighted as both research and supervision: \aprsample{research}{teaching}{research, with supervision}. Chickadee, which I developed to support undergraduate teaching, similarly contributes to both teaching and research: \aprsample{teaching}{research}{teaching, with research}.

\sectionbreak
\cvsection{Significant Research Contributions}
\\ \\
%\textbf{About the figures.} Citation figures are from SciVal (Scopus data, 15 November 2026). Field-weighted citation impact compares a paper's citations with the world average (1.00) for papers of the same field, year and type. Following the University's 2026 position paper on bibliometrics, I give these figures as context for peer judgment, not as measures on their own, and leave out work published in 2026, which has not had time to be cited. Of my \cvnumber{scival-papers} Scopus-indexed papers published up to \cvnumber{scival-last-year}, \cvnumber{scival-top} (\cvnumber{scival-top-share}\%) are among the world's 10\% most-cited for their field, year and type, about twice the expected share.
%\\ \\
\textbf{Making qualitative analysis of online health communities scalable.} Thematic analysis is central to public health research on what people say online, but it does not scale to thousands of posts. My lab built the open-source Computational Thematic Analysis Toolkit (PACM HCI GROUP 2022) and used it to study addiction recovery communities on Reddit (CHI 2022). In a project funded by the Canadian Immunization Research Network (\$105K), we deployed the toolkit with public health researchers and compared their analysis with a manual one (PACM HCI GROUP 2023, published December 2022). The work has since extended to transformer-based topic modelling (ICHEC 2025), to large language models (DeTAILS, CUI 2025 and GI 2026) and to aging in place (PACM HCI 2025). The six papers from 2022--25 have 136 citations. Five are among the world's 10\% most-cited for their field and year: the ICHEC 2025 paper in the top 1\% (from its first 2 citations) and the CHI 2022 paper in the top 2\%. On average the six are cited 6.4 times as often as comparable papers worldwide, or 3.4 times leaving out the ICHEC paper. SciVal places the toolkit papers in a research topic in the top 1\% worldwide for momentum. Trainees from this work now hold positions at Headwaters Health Care Centre, TD and Layer 6.
\\ \\
\textbf{Games for health behaviour change.} My lab designs games and gameful systems that motivate healthier behaviour, and evaluates them with the people they are for. Marcela Bomfim's PhD developed gameful tools for food literacy; her situated shopping app received a Best Paper Honourable Mention at CHI 2020, placing it in the top 5\% of 3,126 submissions. Tina Chan's MSc designed Merlynne, a serious game that motivates peer-to-peer cognitive behavioural therapy (CHI PLAY 2021), and her talk on it won the Best Talk Award at the 2018 Human Factors Inter-University Workshop. Other work addresses stress (CHI PLAY 2021), children with ADHD (2019, 2026) and, in Sebastian Cmentowski's postdoctoral work, sedentary behaviour. Most recently, Veen Wong's PhD (2026, co-supervised with Cayley MacArthur) co-designed Bold Sky with MSW practicum students, a game that engages young men in changing the gender norms behind intimate partner violence (\emph{ACM Transactions on Computing for Healthcare}, 2026; GI 2026). It was presented at the International Summit on Ending Gender-Based Violence (2025) and won the EDI-R Presentation Award at the 2025 Research and Practice in Health Sciences Conference. Related community-based gameful research with Marillac Place supports women and children experiencing homelessness (CHI 2025 poster). The food literacy, Merlynne and stress papers are each among the world's 15\% most-cited for their field and year. I was General Chair of ACM CHI PLAY 2023, and trainees from this work now hold positions at Eindhoven University of Technology (as faculty), Respawn Entertainment and OpenText.
\\ \\
\textbf{Supporting families who care for aging parents at home.} With Leila Homaeian, now a postdoctoral fellow in my lab, and Keiko Katsuragawa, I study how adult children and aging parents share care and stay connected. The program began with a computational analysis of Reddit discussions of informal care for aging parents (PACM HCI 2025), which is cited 2.9 times as often as comparable papers and is among the world's 10\% most-cited for its field and year. It continued with studies in \emph{ACM Transactions on Computing for Healthcare} (2026) and PACM HCI GROUP (2027), and now builds HomeLink, a family communication system (GI 2026 demo). The National Research Council has funded the program twice: \$201K with me as PI (2023--26) and \$321K with me as Co-PI (2026--28). A trainee presenting the work won the Best Oral Presentation Award at the 2024 Research and Practice in Health Sciences Conference.
\\ \\
\textbf{Designing for collaborative sensemaking.} My doctoral and early faculty work studied how groups use shared and personal devices to make sense of information together. It remains my most-cited work. SciVal groups 18 of my papers (2006--22) under interactive tabletop systems for collaboration; together they have 387 citations, and on average they are cited 1.5 times as often as comparable papers worldwide. Three are among the world's 10\% most-cited for their field and year: the NiCE Discussion Room (CHI 2010, top 3\%), collaborative sensemaking on a tabletop and tablets (CHI 2013, top 3\%), and teamwork and taskwork in multi-display groupware (\emph{Personal and Ubiquitous Computing}, 2009, top 7\%). A review of 25 years of empirical research at CSCW received a Best Paper Honourable Mention (CSCW 2018), and I was Program Chair of ACM ISS 2022. The work also reached the public: \emph{Quantum Cats} (2015), a game my lab designed with the Games Institute and the Institute for Quantum Computing (ISS 2016), was featured in IQC's touring \emph{QUANTUM: The Exhibition} and remains an IQC classroom resource; IQC released a browser version in 2026. The same question, how people build shared understanding from information, now drives my lab's tools for qualitative health research.
\\ \\
% TODO (Jim): name the other 2025 workshops in the series, add the CHI 2027 workshop
% outcome (expected 20 Nov 2026), and IH 2026's submission numbers.
% Attendance below: CHI 2025 workshop organizers' email (9 May 2025); Early Career
% Workshop chairs' welcome email to its 15 participants (25 June 2026); IH 2026, CUI 2025
% and CHI PLAY 2023 from Jim (chat, 2026-11-22).
\textbf{Building the HCI and health research community.} Alongside my own research, I help build the international community where human--computer interaction and health research meet. With an international organizing group, I co-organized Envisioning the Future of Interactive Health at CHI 2025, one of 41 workshops accepted from 173 proposals, which drew 111 participants in person and 75 online, and a CHI 2026 meetup on mentorship in HCI and health research. I was a founding convener of the ACM Conference on Interactive Health, whose first conference, in Porto in July 2026, drew about 230 people in person; there I co-chaired its Early Career Workshop for 15 early-career researchers with Geraldine Fitzpatrick and Jochen Meyer. I am its Technical Program Chair for 2027, responsible for the peer review and selection of its archival papers. Scholarly leadership of this kind is rarely recognized with awards; its evidence is that the venues, programs and mentoring now exist for researchers, including my own students, who work between computing and health.
\fi % end APR version only -------------------------------------
% \\ \\
% Finally, in-line with prioritizing publishing at ACM conferences and journals, I dedicate much of my service time to these venues. I actively serve on program committees and as a reviewer for the premiere conferences in HCI, particularly those published by the Association for Computing Machinery (ACM), to which I maintain a professional membership. For instance, I served as General Chair for ACM CHI PLAY 2023 in Stratford, and will be serving as General Chair for ACM CUI 2025 in Waterloo.







% \newpage
% \textbf{PUBLICATIONS:}
% \\ \\
% I typically publish within the Association for Computing Machinery (ACM) research community, and this community has decided that peer-reviewed conference proceedings are the primary venue by which scientific results should be disseminated. In recent years, in recognition that this model is different than other academic traditions, some ACM conferences have begun publishing journals under the ``Proceedings of the ACM in …'' (PACM) series. Nearly all of my publications are now officially therefore journal papers but are associated with ACM conferences. It’s more confusing than it should be.
% \\ \\
% And so, when evaluating my performance \textbf{equal weight should be given to peer-reviewed conference and journal papers}, particularly for those at top-ranked conferences such as the ACM Conference on Human Factors in Computing Systems (CHI), and PACM conference/journal hybrids like the ACM Conference on Computer Supported Cooperative Work (CSCW), the ACM Conference on Computer-Human Interaction in Play (CHI PLAY), and the ACM Conference on Supporting Group Work (GROUP). The review process at these conferences is exceptionally rigorous: full papers (not abstracts) are submitted for anonymous review by 3 expert peers, summarized by 2 associate chairs, and discussed by a program committee that accepts only the most highly rated submissions. These committees place an emphasis on novel research, motivating research for a general audience, and rigour in methods. The selection process is highly competitive, often with acceptance rates approaching 20\%. For full details, please see the ACM Policy on Pre-Publication Evaluation at \href{https://www.acm.org/publications/policies/pre-publication-evaluation}{https://www.acm.org/publications/policies/pre-publication-evaluation}
% \\ \\
% Quantifying the relative importance of these venues within the field of human-computer interaction can be difficult, however Google Scholar provides rankings that can be of some assistance. For example, the impact factor of publications at CHI ($h$-index 122) tends to be much higher than that of journal publications in the field; however, PACM HCI is now ranked 2nd ($h$-index 71) and is placed higher than many top-tier HCI conferences. I would also urge caution when comparing bibliometrics to other fields, including Engineering. \textbf{The ACM’s highly selective review process places an emphasis on quality over quantity}, and thus $h$-indices and other quantitative measures of impact are not directly comparable to other fields. 
% \\ \\
% Further, HCI research often involves many stages, including the design and development of prototype systems, field or laboratory studies of prototype systems, and qualitative and quantitative analysis of collected data. The breadth of tasks often means that many co-authors are listed on a publication, and that each has performed a substantial amount of work on the project based on their own expertise. \textbf{Typically, the first listed author is reserved for the primary researcher (e.g., PhD student) and the last listed author is the supervising researcher or project leader}. My lab has also released software based on the outcomes of our research projects as knowledge mobilization activities. 
% \\ \\
% Finally, in-line with prioritizing publishing at ACM conferences and journals, I dedicate much of my service time to these venues. I actively serve on program committees and as a reviewer for the premiere conferences in HCI, particularly those published by the Association for Computing Machinery (ACM), to which I maintain a professional membership. For instance, I served as General Chair for ACM CHI PLAY 2023 in Stratford, and will be serving as General Chair for ACM CUI 2025 in Waterloo.

\sectionbreak
\cvsection{PUBLICATIONS:}
\\
Google Scholar (22 November 2026): 2{,}211 citations (1{,}242 since 2021); \emph{h}-index 29 (20 since 2021); i10-index 50 (42 since 2021).
\ifAPR
Citation counts depend on the database and the field. Google Scholar indexes the ACM conference proceedings where most HCI research appears, which Scopus and Web of Science cover less fully: a computer scientist with an \emph{h}-index of about 10 in Web of Science can have 20 to 30 in Google Scholar, and \emph{h}-indices differ so much between fields that they should not be compared across them (Hicks et al., 2015, \emph{Nature} 520; University of Waterloo, \emph{Measuring Research Outputs Through Bibliometrics}, 2026). For a comparison within the field, the lines under individual papers give their SciVal field-weighted citation percentiles. A paper's field-weighted citation impact compares its citations with the world average for papers of the same field, year and type (1.00 is that average); its percentile places it among those papers. These figures measure attention from other researchers, not every aspect of a paper's quality.
\fi
\\ \\
% Journal articles by field (data/venues.yaml; regenerated by scripts/venue_profile.py).
% Generated by scripts/venue_profile.py from data/venues.yaml and scholar.bib; do not edit.
Of the 24 refereed journal articles below, 9 are in health journals (\emph{ACM Transactions on Computing for Healthcare}, \emph{BMC Medical Informatics and Decision Making}, \emph{BMJ Innovations}, \emph{Human Factors in Healthcare}, \emph{JMIR Formative Research}, \emph{JMIR Public Health and Surveillance} and \emph{Social Science \& Medicine}), 13 in HCI journals and 2 in a multidisciplinary journal (\emph{PLOS ONE}).
\ifAPR Of the 5 in the review period (2025--2026), 3 are in health journals.\fi

% The APR explains HCI's venues in "A Note on Evaluating Human-Computer Interaction
% Research"; the plain CV carries this short version (wording from that note).
\ifAPR\else
In human--computer interaction, the primary archival venues for full papers are selective ACM conferences, which function as top journals do in other fields; many now publish their papers as journal articles in the Proceedings of the ACM (PACM). Acceptance rates are given under the entries from the most selective venues and tracks (below 30\%).
\fi

\begin{cvparts}

    \nocite{*}

    \cvpart{Refereed Journals:}
    \printcategory[J]{journal}

    \keepwithnext
    \cvpart{Refereed Conference Proceedings (Full and Short Papers):}
    \printcategory[C]{conference}

    \keepwithnext
    \cvpart{Refereed Conference Proceedings (Extended Abstracts, Posters, Demos and Workshops):}
    \printcategory[EA]{abstracts}

    % \cvpart{Workshop Organization:}
    % \printcategory[W]{workshops}

    \keepwithnext
    \cvpart{Technical and Consulting Reports:}
    \printcategory[T]{technicalreports}

    \keepwithnext
    \cvpart{Other Publications:}
    \printcategory[O]{other}

    % \cvpart{Theses and Dissertations:}
    % \printcategory[D]{thesis}

    \keepwithnext
    \cvpart{Patents:}
    \printcategory[P]{patent}

    \sloppy
    \keepwithnext
    \cvpart{Software:}
    \printcategory[S]{software}

\end{cvparts}


\sectionbreak
\cvsection{PRESENTATION TO SCHOLARLY GROUPS:}\\
Speaker(s) denoted by \dag

\begin{cvparts}

    \nocite{*}

    \cvpart*{Invited}
    \printcategory{invited}


    \keepwithnext
    \cvpart*{Self-initiated}
    \printcategory{presentations}

    \keepwithnext
    \cvpart*{Media and Public Engagement}
    \printcategory{media}

\end{cvparts}


\ifAPR % APR version only ------------------------------------------
\sectionbreak
\cvsection{A Note on Research Funding in Human–Computer Interaction}
\\ \\
Human–Computer Interaction (HCI) research draws primarily on funding from the Natural Sciences and Engineering Research Council of Canada (NSERC), particularly its Discovery Grant program. Discovery Grants support independent research programs through discipline-specific peer review of the \textbf{originality, significance and impact} of an investigator's contributions, as well as their training of highly qualified personnel \extlink{https://www.nserc-crsng.gc.ca/NSERC-CRSNG/Reviewers-Examinateurs/DG_eval_crit-eng.asp}{(NSERC Discovery Grant Evaluation Criteria)}. In computing, holding one is the primary indicator of an independently funded, nationally competitive research program.
\\ \\
\textbf{Grant size differs by agency and discipline.} A Discovery Grant typically provides about \$180{,}000 to \$250{,}000 over five years \extlink{https://www.canada.ca/en/research-coordinating-committee/corporate/annual-progress-summary/annex-2.html}{(Tri-Agency Funding Statistics)}. Laboratory disciplines within NSERC receive larger awards that reflect higher research costs \extlink{https://www.nserc-crsng.gc.ca/NSERC-CRSNG/Committees-Comites/EvaluationGroups-GroupesEvaluation_eng.asp}{(NSERC Evaluation Groups and Discipline Structure)}. CIHR Project Grants are typically in the \$850{,}000 to \$980{,}000 range \extlink{https://cihr-irsc.gc.ca/e/54594.html}{(CIHR Project Grant Results)}. These differences reflect the distinct research activities and funding models across disciplines. The SPHS performance review guidelines appropriately recognize that funding should be considered in relation to the nature, scale and objectives of an investigator's research program.
\\ \\
\textbf{My research program has attracted sustained national competitive funding.} I have held an NSERC Discovery Grant continuously since 2015, more than doubling its value at renewal, from \$90{,}000 (2015--22) to \$246{,}000 (2022--28, including a one-year extension), within the typical range above. The National Research Council has funded my aging-in-place research twice: \$201{,}300 with me as PI (2023--26) and \$320{,}966 with me as Co-PI (2026--28).
\\ \\
% The sentence after the lead is generated from grants.bib and data/funders.yaml
% (scripts/funder_profile.py, run by scripts/generate.py).
\textbf{My research also attracts support from health research agencies, reflecting its relevance to public health priorities and interdisciplinary research.} % Generated by scripts/funder_profile.py from data/funders.yaml and grants.bib; do not edit.
I was PI of a Public Health Agency of Canada project (2017) and Co-PI of a Canadian Public Health Association project (2017--18), and a co-investigator or co-applicant on three CIHR grants (2020--21, 2022--24 and 2025--26) and a Canadian Immunization Research Network project (2021--22).
\\ \\
Grants active during the review period are highlighted as \aprsample{research}{}{research}.
\fi % end APR version only -------------------------------------









\sectionbreak
\cvsection{RESEARCH GRANTS AND CONTRACTS:}

\printgrants

% Applications awaiting a decision: grants.bib entries with status = {Submitted}
% (the SPHS guidelines allow listing a grant "in submission"). Details still to come
% print as blank cells (template v0.9.1).
\keepwithnext
\cvsection{RESEARCH GRANTS AND CONTRACTS SUBMITTED:}
\printgrants[Submitted]

\sectionbreak
\ifAPR % APR version only ------------------------------------------
\cvsection{A Note on Graduate Supervision in Human–Computer Interaction}
\\ \\
Human–Computer Interaction integrates methods and perspectives from computer science, health sciences, engineering and the social sciences. My graduate research group reflects this interdisciplinary character, bringing together students from several departments and academic backgrounds. \textbf{My cross-appointment to the Cheriton School of Computer Science further supports this collaborative research and training environment.} Students in my lab work across disciplinary boundaries, developing expertise in computational methods, software practices, human-centred design and digital health technologies. This environment enables health researchers to engage directly with computing and engineering approaches while giving computing students experience addressing substantive health challenges. It advances the University's ``One Waterloo'' emphasis on interdisciplinary collaboration and training.
\\ \\
My approach to graduate supervision emphasizes research participation, mentorship, professional development and timely degree completion. The following summarizes the outcomes of this approach during 2025--26.
\\ \\
\textbf{Funding.} \textbf{It is my policy that every graduate student I supervise is fully funded.} Of the ten graduate students I supervised in 2025--26, nine were funded through my research grants, together with departmental teaching assistantships and, for students I co-supervise, my co-supervisors' grants; the tenth, Veen Wong, held an NSERC Canada Graduate Scholarship.
\\ \\
\textbf{Publishing and presenting.} Trainees were first authors on \cvnumber{review-works-trainee-first} of the \cvnumber{review-works} papers and presentations I published or delivered during 2025--26: \cvnumber{review-journal-trainee-first} of \cvnumber{review-journal} journal articles, \cvnumber{review-fullpapers-trainee-first} of \cvnumber{review-fullpapers} full conference papers, \cvnumber{review-abstracts-trainee-first} of \cvnumber{review-abstracts} extended abstracts, posters, demos and workshops, and \cvnumber{review-presentations-trainee-first} of \cvnumber{review-presentations} presentations. Their work is also among my most cited: trainees first-authored \cvnumber{scival-top-trainee-first} of the \cvnumber{scival-top} papers I published up to \cvnumber{scival-last-year} that are among the world's 10\% most-cited for their field and year, including Robert Gauthier's CHI 2022 paper, in the top 2\%. Veen Wong's presentation of Bold Sky won the EDI-R Presentation Award at the 2025 Research and Practice in Health Sciences Conference.
\\ \\
\textbf{Completion and outcomes.} In 2025--26, Veen Wong completed her PhD in five years (co-supervised with Cayley MacArthur), within the usual time to completion for doctoral students in the School. Ansh Sharma, Karthik Prasad (co-supervised with Oliver Schneider) and Anna Zhelizniak completed their MMath degrees within the program's limits, in five or six terms, and went on to roles at Layer 6 (TD), Tessa AI and ZenduIT. Earlier trainees hold faculty positions at American University and Eindhoven University of Technology, and roles at Headwaters Health Care Centre, Respawn Entertainment and OpenText. I currently supervise four PhD students, two MMath students and a postdoctoral fellow.
\\ \\
\textbf{Supervision approach.} I meet individually with each student weekly and hold a weekly meeting for the entire lab. Students also participate in the larger HCI group meeting, providing opportunities for interdisciplinary exchange and peer mentorship.
\\ \\
\textbf{Committees and mentoring.} I served on 11 graduate committees in 2025 and 12 in 2026, including chairing an oral comprehensive examination. I also examined a thesis at the University of Cape Town, am internal-external examiner for a Waterloo PhD defence in December 2026, and in Fall 2026 mentored a team of three undergraduates in the CRA's UR2PhD program, with my PhD student An Bella Chen as their graduate mentor.
\\ \\
\textbf{Recognition.} In November 2026 my students nominated me for the University of Waterloo Award of Excellence in Graduate Supervision.
\\ \\
Supervision during the review period is highlighted below as \aprsample{teaching}{}{teaching and supervision}. It also appears in the research sections: papers and presentations by students I supervise are shown as \aprsample{research}{teaching}{research, with supervision}, because student-led publications and presentations are central outcomes of HCI research training. External examining is shown as \aprsample{teaching}{service}{supervision, with service}.
\\ \\
\fi % end APR version only -------------------------------------
\keepwithnext[6]
\cvsection{GRADUATE STUDENT SUPERVISION}

\begin{cvparts}

    \cvpart{As Supervisor}

    \printsupervision{PDF}
    


    \printsupervision{PhD}

    \printsupervision{MSc}

    \printsupervision{MMath}

    \printsupervision{MHI}

    \printsupervision{MSW}

    \keepwithnext[13]
    \cvpart{As Committee Member}
    % committees.bib: each committee counts in every year of its date
    \printcommitteecounts{MSc, PhD}

    \ifAPR % APR version only: the named committees ----------------
    \cvsubheading{Committees in the review period:}
    \begin{cvgroup}{PhD}
        \printcommittees[period]{PhD}
    \end{cvgroup}
    \begin{cvgroup}{Master's}
        \printcommittees[period]{MSc}
    \end{cvgroup}
    \fi % end APR version only -----------------------------------

    \keepwithnext[5]
    \cvpart{External Examiner (3)}
    \begin{cvlist}
        \cvitem \hlboth[2025]{teaching}{service}{Sarah Dsane (University of Cape Town, 2025)}
        \cvitem Felwah Alqahtani (Dalhousie, 2022)
        \cvitem Khalid Tearo (Dalhousie, 2021)
    \end{cvlist}

    \keepwithnext[3]
    \cvpart{Undergraduate Research Mentoring}
    \begin{cvlist}
        \cvitem \hlteaching[2026]{CRA Undergraduate Research to PhD (UR2PhD) Mentoring Program, Faculty Mentor, Fall 2026 (team of 3 undergraduate students, with graduate mentor An Bella Chen)}
        \cvitem \hlteaching[2025-2026]{2025 -- 2026: Esraa Hamdy Hamada and Pearl Wiafe (BSW practicum students, Renison School of Social Work; field instructor Veen Wong)}
        \cvitem \hlteaching[2025-]{Spring 2025 -- Present: Maham Kamal (Undergraduate Research Assistant)}
        \cvitem \hlteaching[2025-2026]{2025 -- 2026: Nadeen Al-Janabi (undergraduate honours thesis, HLTH 432A/B)}
        \cvitem \hlteaching[2024]{Fall 2024: Mingchung Xia (Undergraduate Research Assistant)}
        \cvitem \hlteaching[2024]{Spring 2024: Wendy Xu, Ali Raza Zaidi, Abdullah Shahid, Rena Yang and Anna Zhelizniak (Undergraduate Research Assistants); Michelle Ma (co-op, co-supervised with Keiko Katsuragawa, NRC)}
        \cvitem \hlteaching[2024]{Winter 2024: George Wang, Abhinav Jain and Mingchung Xia (Undergraduate Research Assistants)}
        \cvitem \hlteaching[2023]{Fall 2023: Henry Tian, Jean Nordmann and Ryan Lam (Undergraduate Research Assistants)}
        \cvitem \hlteaching[2023]{Fall 2023: Vanessa Duong (USRA, co-supervised with Keiko Katsuragawa, NRC)}
        \cvitem \hlteaching[2023]{Spring 2023: Nicole Mathis, Adrian Davila, Jason Zhao and Peter Li (Undergraduate Research Assistants)}
        \cvitem \hlteaching[2023]{Winter 2023: Yifei Zhang and Tianhao Yao (Undergraduate Research Assistants)}
        \cvitem \hlteaching[2022]{Fall 2022: Atharva Gangal (Undergraduate Research Assistant)}
        \cvitem \hlteaching[2021]{Spring 2021: Talia Wong and Erin Wong (Undergraduate Research Assistants)}
        \cvitem \hlteaching[2020]{Fall 2020: Claudia De Fazio (Undergraduate Research Assistant)}
        \cvitem \hlteaching[2019]{Fall 2019: Nicholas Sia and Kritika Mehta (Undergraduate Research Assistants)}
        \cvitem \hlteaching[2019]{2019: Thomas Huijbregts (undergraduate research assistant)}
        \cvitem \hlteaching[2017-2018]{2017 -- 2018: Colin Whaley (undergraduate honours thesis, HLTH 432A/B)}
        \cvitem \hlteaching[2016-2017]{Fall 2016 -- Winter 2017: Shail Patel (undergraduate thesis)}
        \cvitem \hlteaching[2015-2016]{2015 -- 2016: Adam Baker (Bachelor of Independent Studies honours thesis; one of his supervisors)}
    \end{cvlist}

    \keepwithnext[3]
    \cvpart{Other Mentorship}
    % @service records with level mentorship (activities.bib, service.bib)
    \printservice{mentorship}

    \keepwithnext[3]
    \cvpart{Trainee Awards}
    % awards.bib records with level = {trainee}: awards my students won for
    % work I supervised (also in apralsoteaching)
    \printawards[trainee]
\end{cvparts}

\sectionbreak
\ifAPR % APR version only ------------------------------------------
\cvsection{A Note on Service in Human–Computer Interaction}
\\ \\
Human–Computer Interaction (HCI) is an international research community whose scholarly activities depend on extensive collaboration and professional service. The field's leading ACM conferences bring together international networks of reviewers, program committee members, associate chairs and conference leaders to evaluate research and advance the discipline. \textbf{Conference leadership is an important form of scholarly contribution.} General Chairs provide strategic and organizational leadership, while Program Chairs oversee peer review and the selection of archival research publications. Together, these roles sustain the international scholarly infrastructure through which researchers share, evaluate and advance knowledge.
\\ \\
My service during the review period has focused on building the HCI and health research community while contributing to University and Faculty governance. As General Chair of ACM CUI 2025, I brought an international ACM conference of about 130 attendees to Waterloo. As a founding convener of the ACM Conference on Interactive Health, I helped establish a new international venue connecting researchers in computing and the health sciences; it launched in 2026 with about 230 attendees in person, and I am its Technical Program Chair in 2027. Within the University, I served as Chair Elect of the Health Faculty Council in 2025--26 and now serve as its Chair in 2026--27, contributing to Faculty-wide academic governance. Earlier, I was General Chair of ACM CHI PLAY 2023 in Stratford (about 140 attendees in person) and Program Chair of ACM ISS 2022, and I review regularly for leading HCI conferences and journals.
\\ \\
Service during the review period is highlighted as \aprsample{service}{}{service}. Peer review, program committees, editorial responsibilities and conference leadership also demonstrate scholarly standing and contributions to the research community. They are therefore identified as service with research: \aprsample{service}{research}{service, with research}. Leadership roles appear under both SCHOLARLY AND PROFESSIONAL ACTIVITIES and SERVICE.
\\ \\
\fi % end APR version only -------------------------------------
\keepwithnext
\cvsection{SERVICE}

% the Faculty of Health's parts, in its order, with plain titles
\begin{cvparts}
    \cvpart*{University}
    \printservice{university}

    \cvpart*{Faculty}
    \printservice{faculty}

    \cvpart*{Department}
    \printservice{school}

    \cvpart*{Community}
    \printservice{discipline}
\end{cvparts}

\ifAPR % APR version only ------------------------------------------
\sectionbreak
\cvsection{A Note on Teaching Innovation: Chickadee}
\\ \\
Chickadee is a programming and assessment platform I designed and built to improve how students learn through programming. It provides immediate feedback on submitted code, allowing students to identify and correct mistakes while course material is still fresh and to learn through repeated practice. Students submit their work in a web browser, in Python, R or five other languages; Chickadee runs the instructor's tests on each submission, and students can resubmit as often as they need. It replaces Marmoset, an aging grading system from the University of Maryland, and is hosted by IST at chickadee.uwaterloo.ca, uses the University's single sign-on, sends grades to LEARN, and is open source \extlink{https://github.com/JimWallace/Chickadee}{(github.com/JimWallace/Chickadee)}.
\\ \\
\textbf{Inclusion.} Chickadee supports accessible participation: students use it entirely in the browser, with no software installation required, so all students can take part from the first week of classes, without one-on-one technical support for each student's computer. Unlimited attempts with immediate feedback let every student practise until their work is right. Chickadee also has built-in support for slip days, and automatically (and continuously) synchronizes grades with LEARN.
\\ \\
\textbf{Student benefit.} Chickadee's first deployment was in HLTH 230 (Introduction to Health Informatics) in Spring 2026, providing an opportunity to evaluate its use in a large undergraduate course; Winter 2026 ran on Marmoset. Its 143 students made 8,904 submissions across 17 assignments, with a median wait for results of 7.5 seconds, and 62\% of students' assignments were resubmitted at least once after the first results. In the Student Course Perceptions survey, the course scored 4.2 out of 5 overall and 4.5 for creating a supportive environment for learning, and students named what Chickadee provides: ``being able to get immediate feedback on assessments'' and ``unlimited attempts to complete the assignments/labs.''
Chickadee has attracted interest beyond my own teaching. SPHS colleagues plan to adopt the platform, and a Teaching Innovation Incubator proposal in preparation would extend its use to seven courses and more than 1,300 students across Health and Engineering in 2027--28. Its open-source release and the WatITis talk make this major instructional innovation available beyond my own courses. Colleagues reviewed my teaching in 2024 (HLTH 454/629 and HLTH 637) as part of my promotion review, and I plan a peer review of teaching with Chickadee in 2027.
\\ \\
\textbf{Institutional infrastructure and collaboration.} It runs on University servers through IST, uses University single sign-on, and synchronizes grades with LEARN. It passed the University's Information Risk Assessment and Privacy Impact Assessment under Policy 46, including independent security testing. Grading runs in isolated, secured environments, and lab notebooks run in the student's browser. I developed Chickadee in collaboration with colleagues in IST, the Privacy Office and the Faculty of Health. The project demonstrates how software development, institutional IT expertise, privacy and security practices, and pedagogical knowledge can be brought together to create practical infrastructure for teaching. I will present it at WatITis 2026, the University's IT conference, on December 9.
%\\ \\
% Chickadee and the WatITis talk appear under PUBLICATIONS and PRESENTATION TO SCHOLARLY GROUPS, shown as \aprsample{teaching}{research}{teaching, with research}.

\sectionbreak
\cvsection{A Note on Teaching: Grades in HLTH 230}
\\ \\
I redeveloped HLTH 230 while teaching it in Winter 2026 and further refined it for Spring in response to student feedback. The revisions strengthened the course's framing around health applications, and Spring 2026 introduced Chickadee for all 17 programming assignments. Student Course Perceptions results improved across all six teaching items between the two offerings. In particular, stimulating interest increased from 3.3 to 4.0, identifying intended learning outcomes increased from 3.7 to 4.3, and assessing those outcomes through graded work increased from 3.6 to 4.2. Student achievement remained consistent with previous offerings. Mean final grades were 84.4\% in Winter 2026 and 82.9\% in Spring 2026, compared with 84.6\% in Winter 2020, my previous offering, and within the range of the course's other offerings since 2017 (77.2\% to 97.3\%), as the figure below shows.
\\ \\
% HLTH 230 final course grades by term (APR version only), drawn from
% data/hlth230-grades.csv (see data/hlth230-grades.md). Shaded: terms Jim taught.
\begin{center}
\begin{tikzpicture}
\begin{axis}[
    width=\linewidth, height=6.2cm,
    ymin=60, ymax=104, ytick={60,70,80,90,100},
    ylabel={Final grade (\%)}, ylabel style={font=\small},
    xmin=-0.6, xmax=17.6,
    xtick=data, xticklabels from table={data/hlth230-grades.csv}{term},
    xticklabel style={font=\scriptsize, rotate=60, anchor=east},
    yticklabel style={font=\scriptsize},
    ymajorgrids, grid style={black!12},
    unbounded coords=jump,
    legend style={font=\scriptsize, draw=none, fill=none, at={(0.5,1.02)}, anchor=south},
    legend columns=3,
    table/col sep=comma,
  ]
  % terms I taught: Winter 2020 (4), Winter 2026 (16), Spring 2026 (17)
  \fill[black!8] (axis cs:3.5,60) rectangle (axis cs:4.5,104);
  \fill[black!8] (axis cs:15.5,60) rectangle (axis cs:17.5,104);
  \addplot[blue!70!black, thick, mark=*, mark size=1.4pt]
    table[x expr=\coordindex, y=mean] {data/hlth230-grades.csv};
  \addplot[green!50!black, thick, mark=*, mark size=1.4pt]
    table[x expr=\coordindex, y=median] {data/hlth230-grades.csv};
  \addplot[orange!85!black, thick, dashed, mark=*, mark size=1.4pt]
    table[x expr=\coordindex, y=mode] {data/hlth230-grades.csv};
  \legend{Mean, Median, Mode}
\end{axis}
\end{tikzpicture}

\parbox{\linewidth}{\footnotesize HLTH 230 final course grades by term, Winter 2017 -- Spring 2026. Shaded: terms I taught. Winter 2017 -- Spring 2025 from the University's Institutional Analysis and Planning count-date data (Quest); Winter and Spring 2026 from LEARN final grades, where the mode is not comparable and is not shown.}
\end{center}

\fi % end APR version only -------------------------------------

\sectionbreak
\keepwithnext[6]
\cvsection{COURSES TAUGHT}

% courses.bib; Student Course Perceptions scores in the APR version only
\renewcommand\cvcoursescplabel{SCP}
\renewcommand\cvcourseaveragelabel{SPHS Average}
\ifAPR\printcourses[scp]\else\printcourses\fi

\keepwithnext
\cvsection{PROFESSIONAL DEVELOPMENT AND TRAINING}

% service.bib records with level = {training}
\printservice{training}

\sectionbreak[6\baselineskip]
\cvsection{AREAS OF TEACHING EXPERTISE} \\
Health Informatics, Digital Health, Human-Computer Interaction, Computational Social Science, Games, Computer-Supported Cooperative Work, Information Visualization, Machine Learning, Artificial Intelligence
\\ \\
\cvsection{CURRENT RESEARCH INTERESTS} \\
I design, build and evaluate interactive technologies for health, working with the people who will use them.
\\ \\
\textbf{Games for health behaviour change:} games and gameful systems that motivate healthier behaviour, co-designed and evaluated with the people they are for, including food literacy, peer support for mental health, stress, children with ADHD, sedentary behaviour, and gender-transformative games that engage young men in preventing intimate partner violence.
\\ \\
\textbf{Human-centred qualitative data science for public health:} computational and large-language-model support for thematic analysis of online health communities, in tools that keep researchers in control of interpretation.
\\ \\
\textbf{Technology for aging in place and family care:} communication and awareness technologies that help aging parents and the adult children who care for them share care and stay connected.
\\ \\
\textbf{Collaborative sensemaking:} how groups build shared understanding from information with shared and personal devices, the foundation of my lab's tools for qualitative health research.


% The letter on process in SPHS reviews (letters/process.tex; external-input-letter.tex on
% its own) is not in the APR packet yet (Jim, 2026-11-24). To add it after the CV:
% \ifAPR% Letter to the Director and the APR Committee on process in the School's annual performance
% reviews, on SPHS letterhead, signed by Jim and enclosed with the APR (after cover-letter.tex).
% It began as the CV's "A Note on External Input for Interdisciplinary Work" and was widened to
% several process issues for the Committee to discuss (Jim, 2026-11-24). It documents process,
% not individuals: it names no one, refers to no individual review, and quotes only public
% instruments. Compile on its own (Overleaf: Menu > Main document > external-input-letter.tex).
%
% Sources: Memorandum of Agreement 13.5.1(d), 13.5.3, 13.5.6(a) and (c), 13.5.7, 13.5.10(a), 13.3.3(d)
% (texts from Jim, chat, 2026-11-24, and his own documents); SPHS Addendum to the Faculty of
% Health Performance Evaluation Guidelines (2025--26): "Research and Scholarship",
% "Recognition of multidisciplinarity", "Approach to performance assessment", "Use of
% collective indicators" and "Importance of Feedback" (wording as Jim quotes it); SPHS
% Constitution s. 3.1.6 as amended in November 2026 (School vote results, 13 November 2026);
% Faculty of Health Performance Review Guidelines, "Anomaly Adjustments" (from the version
% circulated with the revisions proposed for Fall 2026; Jim, 2026-11-24);
% Report of the FRC Working Group on Faculty Performance (8 January 2026); CRA, "Hiring and
% Promotion Guidance for Computing Researchers" (September 2025); Wallace (2026),
% doi:10.5281/zenodo.24397928 (findings as its README states them).
% Research on groups: the five references at the end; each record checked against Crossref
% (2026-11-24); claims from the abstracts where available (Stasser & Titus 1985; Latane 1981;
% Lu et al. 2012); Stasser, Taylor & Hanna 1989 and Noelle-Neumann 1974 as commonly summarized.
%
% The text of the letter. The standalone document (external-input-letter.tex) and the APR version
% of the CV (main.tex, \ifAPR) both read this file; letters/aprletter.tex defines the
% aprletter environment for each.
\begin{aprletter}{%
The Director and the Annual Performance Review Committee\\
School of Public Health Sciences\\
University of Waterloo}{Dear Interim Director and Members of the APR Committee,}

\textbf{Re: Process in the School's annual performance reviews}

I write to ask that the Committee discuss, before it begins this cycle's reviews, six features of how the School conducts annual performance reviews. For each, I note what the Memorandum of Agreement (MoA) or the SPHS Addendum provides and propose a step the School could take. I also ask, for my own file, that expert opinion be obtained under the Addendum (item~2).

\textbf{1. The size of the committee.} The MoA provides that, in a department of more than 15 full-time equivalent regular faculty positions, the Members ``shall elect an advisory committee of no more than five Members to assist the Chair'' (MoA 13.5.6(c)). School Council has amended Section 3.1.6 of the School Constitution to provide a committee of five, ``as per the MoA 13.5.6(c),'' but the committee for this cycle was elected with ten. The limit matters for more than form. The influence a group exerts on an individual grows with the number of people exerting it (Lat\'an\'e, 1981), so in a larger committee the norms most of its members share carry more weight; members who perceive their view to be in the minority are less likely to voice it (Noelle-Neumann, 1974); and groups discuss the information their members hold in common and neglect information that only one member holds (Stasser \& Titus, 1985), an effect a meta-analysis of 65 studies confirms (Lu, Yuan, \& McLeod, 2012) and that is stronger in larger groups (Stasser, Taylor, \& Hanna, 1989). In a performance review, the member who knows the norms of a minority discipline is the one holding that unshared information. A larger committee therefore makes it more likely that the norms of the School's majority disciplines are applied to everyone, and less likely that expert opinion, once obtained, is heard. \emph{Proposal:} that the Committee consider how this cycle's review can meet the MoA's limit, and that it record how it will weigh expert opinion when it obtains it.

\textbf{2. Disciplinary expertise and expert opinion.} The Addendum recognizes that the School is multidisciplinary and that ``performance indicators should be interpreted in a manner consistent with these various disciplinary perspectives and traditions''; it notes that ``for some disciplines (e.g., engineering and computer science), peer-reviewed conference proceedings are also regarded as important.'' It provides the means: ``Where assessment of quality and impact cannot be readily evaluated because of inadequate domain expertise, the Director shall have the discretion to obtain expert opinion, in confidence, from faculty members of the School, or elsewhere in the University, who hold regular appointments in a relevant disciplinary area.'' The Addendum places this responsibility on the Director. It does not make expert opinion depend on a member's request: a request may be sufficient to prompt it, but it is not necessary. Nor does the Addendum say when expert opinion will be considered, how an expert is chosen, what the expert receives, or how the committee uses the advice. Without a defined process, whether a member's work is read by someone in their own field depends on whether the member knows that expert opinion can be obtained, and feels able to ask; and a member who does ask is left to collate the evidence, analyze it and plead their case, which is not equitable. The Computing Research Association gives the same advice for computing: when a candidate's field has norms that differ from those of the unit, the unit should document how it will evaluate the candidate against the criteria of that field (\emph{Hiring and Promotion Guidance for Computing Researchers}, 2025). \emph{Proposal:} a defined, equitable and transparent process, added to the Addendum (which MoA 13.5.1(d) requires to be posted and publicly accessible), that
\begin{enumerate}[leftmargin=1.75em, itemsep=2pt, topsep=4pt, parsep=0pt, label=(\alph*)]
    \item has the Director determine, for every member's file and not only on request, whether the committee has the expertise to assess its quality and impact, and obtain expert opinion where it does not, with any member also able to request it;
    \item draws the expert, as the Addendum provides, from faculty with regular appointments in the relevant area, with no conflict of interest with the member;
    \item gives the expert the member's file, the Addendum and the Faculty's rating descriptors;
    \item records in the member's file the Director's decision, whether expert opinion was obtained and, if so, how the committee took it into account.
\end{enumerate}
\emph{My request:} my research is in human--computer interaction (HCI), a computing discipline whose norms for publication, authorship, funding and service differ from those of health informatics and much of public health, and no member of the committee works in HCI. I ask that expert opinion be obtained on my 2025--2026 file from a faculty member who holds a regular appointment in HCI, for example in the Cheriton School of Computer Science, where I am cross-appointed; who publishes regularly in HCI venues, in particular those of ACM SIGCHI; and who can comment on both publication and service in those communities. I can suggest names, and will disclose any relationship I have with each. I also ask that your decision on this request, and your reasons, be recorded in my performance review.

\textbf{3. Ratings on each member's record, against the descriptors.} The MoA gives the Chair ``the responsibility for annual performance evaluations of all Members in the Department'' (MoA 13.5.6(a)) and sets out the ratings and their descriptions (MoA 13.5.3). The Dean ``may modify the ratings for a Member or Members of a Department, if necessary, to maintain consistency of standards across the Faculty,'' and ``shall inform the Chair in writing of the final individual and overall ratings, together with reasons for any changes'' (MoA 13.5.7). Nothing in these provisions sets a unit average or a distribution of ratings. The University's FRC Working Group on Faculty Performance (January 2026) found practices ``lacking clarity on how the full spectrum of the FPR scores available (MoA Article 13.5.3) is to be interpreted and applied.'' \emph{Proposal:} that the Committee confirm that each rating is set on the member's own record against the descriptions in MoA 13.5.3 and the Addendum, using the full scale, and that no unit average or expected distribution is used in setting ratings.

\textbf{4. Feedback that addresses quality and impact.} The Addendum requires that the indicators used ``must have face validity to enable individual faculty members, the review committee, the Director and other stakeholders to assess quality/excellence, innovation/creativity, and impact,'' and asks that written feedback be ``honest, concise, and constructive,'' noting that ``it may be useful to include specific examples of excellence or areas of challenge.'' Feedback that only counts outputs would not meet this standard, and would not show how disciplinary norms (item~2) were applied. \emph{Proposal:} that comments on the evaluation form address quality, innovation and impact in each area, with examples, and say how the norms of the member's discipline were taken into account.

\textbf{5. Disagreements.} The MoA provides that ``a Member who disagrees with her/his performance evaluation should proceed first to the Department Chair, and then, if not resolved, to the Dean of the Faculty for disposition'' (MoA 13.5.10(a)). It does not set a timeline, and the School has no written procedure for it. Where the MoA provides for changes to ratings, it requires written reasons (MoA 13.5.7). \emph{Proposal:} a written procedure, with a timeline that allows a disagreement to be resolved within the cycle, and a written response with reasons at each step.

\textbf{6. Salary anomalies.} The MoA establishes an annual Anomalies Fund in each Faculty ``to correct individual salary anomalies'' (MoA 13.3.3(d)), and the Faculty of Health Performance Review Guidelines provide that these adjustments ``are determined by the Dean based upon recommendations of the Chair/Directors.'' Identifying anomalies in the School is therefore part of the Director's role, but the School has no written procedure for it. In practice a member must find an anomaly in the public salary disclosures, analyze it and make the case, the same burden as in item~2, and one that falls hardest on members who already feel least able to ask. \emph{Proposal:} that the Director review the salaries of all members each year, by rank and by years since appointment or promotion, identify possible anomalies and recommend adjustments to the Dean, and report to the School, without names, how many were identified and recommended.

\textbf{Why these matter.} Each item matters most to members whose work follows the norms of another discipline, interdisciplinary and cross-appointed faculty most of all, who are the most likely to be assessed against criteria that do not fit their field. Each proposal also moves the work of making the case from the individual member, who would otherwise have to collate the evidence, analyze it and plead for its fair reading, to the process itself, where it belongs. Salary is one place where the effect would show, because merit increases accumulate in it. My 2026 analysis of the University's public salary disclosures for SPHS from 2011 to 2025 (doi:10.5281/zenodo.24397928) finds that about 80\% of the person-years of faculty affiliated with the Master of Health Informatics program fall below the School's same-year median salary, against about 41\% for other SPHS faculty, and estimates their salary growth at about \$1,595 a year slower. The analysis is descriptive: public data cannot identify the mechanism behind the gap, and the report recommends a faculty-level equity audit that links salary outcomes to rank, workload and performance reviews.

I ask that these items be placed on the Committee's agenda before it begins this cycle's reviews, and I would be glad to help draft the procedures proposed here.

{\small
\raggedright
\par\medskip\noindent\textbf{References}\par\smallskip
\noindent Lat\'an\'e, B. (1981). The psychology of social impact. \emph{American Psychologist}, 36(4), 343--356. \url{https://doi.org/10.1037/0003-066X.36.4.343}\par
\noindent Lu, L., Yuan, Y. C., \& McLeod, P. L. (2012). Twenty-five years of hidden profiles in group decision making: A meta-analysis. \emph{Personality and Social Psychology Review}, 16(1), 54--75. \url{https://doi.org/10.1177/1088868311417243}\par
\noindent Noelle-Neumann, E. (1974). The spiral of silence: A theory of public opinion. \emph{Journal of Communication}, 24(2), 43--51. \url{https://doi.org/10.1111/j.1460-2466.1974.tb00367.x}\par
\noindent Stasser, G., Taylor, L. A., \& Hanna, C. (1989). Information sampling in structured and unstructured discussions of three- and six-person groups. \emph{Journal of Personality and Social Psychology}, 57(1), 67--78. \url{https://doi.org/10.1037/0022-3514.57.1.67}\par
\noindent Stasser, G., \& Titus, W. (1985). Pooling of unshared information in group decision making: Biased information sampling during discussion. \emph{Journal of Personality and Social Psychology}, 48(6), 1467--1478. \url{https://doi.org/10.1037/0022-3514.48.6.1467}\par
}
\aprclosing{Sincerely,}
\end{aprletter}
\fi
\end{document}










